% ======================================================================
%  Preambolo — modello di tesi DISI (cartabinaria/latex-thesis) più i
%  pacchetti e le macro proprie di questo lavoro.
%  https://github.com/cartabinaria/latex-thesis
% ======================================================================

\usepackage[utf8]{inputenc}
\usepackage[T1]{fontenc}
% Gabbia piu' larga del modello, ma il margine INTERNO resta 3,5 cm come da
% modello: e' quello che la rilegatura mangia, e ridurlo si vedrebbe sul lato
% cucito. Si guadagna spazio sull'esterno (2,5 -> 2 cm) e soprattutto in
% verticale, che il modello non fissa affatto e lasciava al default di
% geometry, circa 4,5 cm per lato su A4.
\usepackage[a4paper,inner=3.5cm,outer=2cm,top=2.5cm,bottom=2.5cm]{geometry}

\usepackage[titletoc,title,toc,page]{appendix}
\usepackage{verbatim}
\usepackage{placeins}
\usepackage{listings}
\lstdefinestyle{kairos}{
  basicstyle=\ttfamily\footnotesize,
  keywordstyle=\bfseries,
  commentstyle=\itshape\color{gray!60!black},
  morecomment=[l]{//},
  morekeywords={procedure,local,delocal,call,uncall,if,then,else,fi,from,do,
                loop,until,par,and,rap,try,rollback,yrt,int,stack,channel,
                push,pop,ssend,srecv,show,empty,nil},
  columns=fullflexible, breaklines=true, showstringspaces=false,
  frame=leftline, framesep=6pt, rulecolor=\color{gray!40}, xleftmargin=8pt,
}
\lstset{style=kairos}
\usepackage[italian]{babel}
% Il greco serve per due parole sole (l'etimologia di Kairos). Questa
% installazione di TeX non ha i font LGR ne' babel-greek, quindi le lettere si
% compongono in matematica, che e' sempre disponibile.
\newcommand{\kairosgr}{\ensuremath{\kappa\alpha\iota\rho\acute{o}\varsigma}}
\newcommand{\chronosgr}{\ensuremath{\chi\rho\acute{o}\nu o\varsigma}}
\usepackage{tikz}
\usepackage{parskip}

\usepackage{graphicx}
\usepackage{chngcntr}
\counterwithin{table}{chapter}

\usepackage{fancyhdr}
\usepackage{indentfirst}
\usepackage{float}
\usepackage{soul}
\usepackage[font=footnotesize,labelfont=bf]{caption}

\usepackage{multirow}
\usepackage{hyphenat}
% Solo parole senza lettere accentate: \hyphenation non le accetta con
% questo inputenc.
\hyphenation{mate-mati-ca recu-perare con-cor-ren-za sin-cro-niz-za-zio-ne
             se-man-ti-ca in-ver-sio-ne}

\usepackage{lscape}

% Ultima risorsa di TeX prima di far sporgere una riga nel margine: le concede
% un po' di elasticita' in piu' sugli spazi. Serve dove una formula in linea
% lunga non ha punti di a-capo legali.
\setlength{\emergencystretch}{3em}

% ---- aggiunte a questo lavoro rispetto al modello -----------------------
% Il modello non le carica perché non prevede contenuto matematico.
\usepackage{amsmath,amssymb}
\usepackage{array}          % colonne >{$}l<{$} nel quadro riassuntivo
\usepackage{xcolor}         % \static, \dynamic, \segnaposto
% ------------------------------------------------------------------------

\usepackage{natbib}
\bibliographystyle{alpha}
\setcitestyle{super,open={[},close={]}}

\newcommand{\rom}[1]{\uppercase\expandafter{\romannumeral #1\relax}}

\usepackage{pdfpages}

% hyperref e cleveref vanno per ultimi, in quest'ordine.
\usepackage{hyperref}
\hypersetup{colorlinks=true,urlcolor=blue!55!black,linkcolor=blue!55!black,
            citecolor=blue!55!black}
\usepackage[italian,capitalize,noabbrev]{cleveref}
\usepackage{pgfplots}
\pgfplotsset{compat=1.18}
% ======================================================================
%  Notazione propria
% ======================================================================

% ATTENZIONE -- ordine degli argomenti di \cfg.
% Gli argomenti sono, per ragioni storiche, (comando, store): \cfg{c}{\sigma}.
% L'output reso è invece (store, comando): <sigma, c>, per allinearsi alla
% semantica di base [LamiLaneseStefani2024], che scrive <sigma, s>.
% Il corpo della macro scambia quindi #1 e #2. NON riscrivere i ~100 siti di
% chiamata: si cambia solo qui.
\newcommand{\cfg}[2]{\langle #2,\, #1 \rangle}          % argomenti (c,sigma) -> reso <sigma, c>
\newcommand{\upd}[3]{#1[#2 \mapsto #3]}                 % sigma[v -> n]
\newcommand{\updc}[4]{#1[#2\langle #3\rangle \mapsto #4]}  % sigma[a<k> -> n]: sola cella k
\newcommand{\toe}{\to_e}
\newcommand{\tob}{\to_b}
\newcommand{\toc}{\to_c}
\newcommand{\ttt}{\mathit{tt}}
\newcommand{\ff}{\mathit{ff}}
% L'esito d'errore è \bot, come nella semantica di base [LamiLaneseStefani2024].
% La macro \err è mantenuta (ridefinita) per non toccare i ~50 siti di chiamata.
\newcommand{\err}{\bot}
\newcommand{\xor}{\oplus}
\newcommand{\Inv}[1]{\mathcal{I}(#1)}
\newcommand{\Store}{\mathit{Store}}

% nome regola a destra
\newcommand{\rname}[1]{\quad(\textsc{#1})}
% regola d'errore a runtime (fa parte di ->_c, produce err)
\newcommand{\rnameerr}[1]{\quad(\textcolor{red!60!black}{\textsc{#1}})}
% Vincoli di buona formazione: stanno fuori dalla relazione ->_c, e si
% distinguono per QUANDO li si verifica.
%   \static  (verde)   verificato sul testo del programma, in compilazione
%   \dynamic (arancio) sorvegliato durante l'esecuzione; violarlo da' \err
\newcommand{\static}[1]{\textcolor{green!40!black}{#1}}
\newcommand{\dynamic}[1]{\textcolor{orange!80!black}{#1}}

% ---- segnaposto ----
% Capitoli ancora da scrivere: resi in grigio, titolo e corpo, per essere
% immediatamente riconoscibili come tali.
% Identificatori e frammenti di codice nel testo corrente.
\newcommand{\code}[1]{\texttt{#1}}

\newcommand{\segnaposto}[1]{\textcolor{gray}{#1}}
\newenvironment{segnapostotesto}{\color{gray}\itshape}{}
